% META: {"id": "1SPE-SUITES-COURS-07-FR", "chapitre": "1SPE-SUITES", "type_objet": "cours", "sous_type": "td_fil_rouge", "titre": "TD fil rouge — Deux offres d'epargne", "temps_gabarit": 7, "mode_creation": "ex_nihilo", "sources_inspiration": [], "status": "generated", "capacites_codes": ["C2", "C3", "C6", "C7"]}
% BEGIN-VERIFY
% from sympy import symbols, Rational, solve, simplify
% n = symbols('n', nonneg=True, integer=True)
% # Offre A : suite arithmetique A_n = 200*n
% A = lambda k: 200 * k
% assert A(0) == 0
% assert A(1) == 200
% assert A(12) == 2400
% # Offre B : suite geometrique B_n = 1000 * 1.004^n
% B = lambda k: 1000 * (1.004 ** k)
% assert abs(B(0) - 1000) < 1e-9
% assert abs(B(1) - 1004) < 1e-9
% # B demarre en tete (capital initial 1000 euros) puis A reprend l'avantage des le rang 6
% assert B(5) > A(5)
% assert B(6) <= A(6)
% for j in range(7, 10):
%     assert B(j) < A(j)
% # Recherche du premier depassement de B sur A a partir du rang 10
% k = 10
% while B(k) <= A(k):
%     k += 1
% assert k == 1415
% assert B(1415) > A(1415)
% assert B(1414) <= A(1414)
% assert B(1416) - B(1415) > 200
% assert abs((B(1416) - B(1415)) - 0.004 * B(1415)) < 1e-9
% # Verification valeurs approximatives
% assert abs(B(1415) - 1000 * 1.004**1415) < 1e-6
% assert A(1415) == 283000
% END-VERIFY

\section*{Temps 7 — TD fil rouge : Deux offres d'épargne}

\textbf{Situation.} Un établissement bancaire propose deux formules d'épargne. On note $n$ le nombre de mois écoulés depuis la souscription.

\begin{itemize}
  \item \textbf{Offre A} : versement fixe de $200$~€ par mois sur un livret initialement vide et sans intérêts.
  \item \textbf{Offre B} : capital initial de $1\,000$~€ placé à intérêts composés au taux mensuel de $0{,}4\,\%$.
\end{itemize}

On note $A_n$ le capital de l'offre A et $B_n$ le capital de l'offre B après $n$ mois.

\bigskip

%---------------------------------------------------------------
% PARTIE A — Offre A
%---------------------------------------------------------------

\begin{exercice}{1SPE-SUITES-COURS-07-FR-EX1}{2}{15}

\textbf{Partie A — Offre A : modèle arithmétique}

\begin{enumerate}
  \item Exprimer $A_n$ en fonction de $n$. Quelle est la valeur de $A_0$ ?
  \item Montrer que la suite $(A_n)$ est arithmétique. Préciser son premier terme et sa raison.
  \item Calculer $A_{12}$, $A_{24}$ et $A_{60}$. Interpréter ces valeurs.
\end{enumerate}

\end{exercice}

\begin{corrige}{1SPE-SUITES-COURS-07-FR-EX1}

\commentaireMarge{Chaque mois, le capital augmente du même montant : évolution absolue constante $\Rightarrow$ modèle arithmétique.}

\textbf{Question 1.} Avant tout versement ($n = 0$), le livret est vide : $A_0 = 0$. Chaque mois, on ajoute $200$~€, donc :
\[
  A_n = 200n.
\]

\textbf{Question 2.} On calcule la différence entre deux termes consécutifs :
\[
  A_{n+1} - A_n = 200(n+1) - 200n = 200.
\]
La différence est constante et égale à $200$. Donc $(A_n)$ est une suite arithmétique de premier terme $A_0 = 0$ et de raison $r = 200$.

\textbf{Question 3.}
\[
  A_{12} = 200 \times 12 = 2\,400 \text{ €} \qquad \text{(capital après 1 an),}
\]
\[
  A_{24} = 200 \times 24 = 4\,800 \text{ €} \qquad \text{(capital après 2 ans),}
\]
\[
  A_{60} = 200 \times 60 = 12\,000 \text{ €} \qquad \text{(capital après 5 ans).}
\]

\end{corrige}

\bigskip

%---------------------------------------------------------------
% PARTIE B — Offre B
%---------------------------------------------------------------

\begin{exercice}{1SPE-SUITES-COURS-07-FR-EX2}{2}{15}

\textbf{Partie B — Offre B : modèle géométrique}

\begin{enumerate}
  \item Exprimer $B_{n+1}$ en fonction de $B_n$. En déduire que $(B_n)$ est une suite géométrique. Préciser son premier terme et sa raison.
  \item Exprimer $B_n$ en fonction de $n$.
  \item Calculer $B_{12}$, $B_{24}$ et $B_{60}$ (arrondir à l'euro). Interpréter.
\end{enumerate}

\end{exercice}

\begin{corrige}{1SPE-SUITES-COURS-07-FR-EX2}

\commentaireMarge{Taux mensuel de $0{,}4\,\%$ : coefficient multiplicateur $q = 1 + 0{,}004 = 1{,}004$.}

\textbf{Question 1.} Chaque mois, le capital est multiplié par $q = 1{,}004$ (augmentation de $0{,}4\,\%$) :
\[
  B_{n+1} = 1{,}004 \times B_n.
\]
Le rapport $\dfrac{B_{n+1}}{B_n} = 1{,}004$ est constant et $B_0 = 1\,000
\neq 0$. Donc $(B_n)$ est une suite géométrique de premier terme $B_0 = 1\,000$ et de raison $q = 1{,}004$.

\textbf{Question 2.} Par la formule du terme général d'une suite géométrique :
\[
  B_n = B_0 \times q^n = 1\,000 \times 1{,}004^n.
\]

\textbf{Question 3.}
\[
  B_{12} = 1\,000 \times 1{,}004^{12} \approx 1\,000 \times 1{,}0489 \approx 1\,049 \text{ €,}
\]
\[
  B_{24} = 1\,000 \times 1{,}004^{24} \approx 1\,000 \times 1{,}1005 \approx 1\,101 \text{ €,}
\]
\[
  B_{60} = 1\,000 \times 1{,}004^{60} \approx 1\,000 \times 1{,}2706 \approx 1\,271 \text{ €.}
\]

Sur 5 ans, le solde de l'offre B augmente d'environ $271$~€. La comparaison
avec le solde de l'offre A doit tenir compte de flux de versements différents :
A reçoit $200$~€ chaque mois, tandis que B ne reçoit que le capital initial de
$1\,000$~€.

\end{corrige}

\bigskip

%---------------------------------------------------------------
% PARTIE C — Comparaison algorithmique
%---------------------------------------------------------------

\begin{exercice}{1SPE-SUITES-COURS-07-FR-EX3}{2}{20}

\textbf{Partie C — À partir de quel mois $B_n > A_n$ ?}

\begin{enumerate}
  \item Vérifier que $B_5 > A_5$ et que $B_6 \leqslant A_6$. Que peut-on en conclure sur les tout premiers mois ?
  \item Lire et expliquer le programme Python ci-dessous, qui détermine
  le premier rang $n$ (à partir du rang $10$) tel que $B_n>A_n$.
  \item Exécuter mentalement les deux premiers tours de boucle (vérifier les valeurs de \texttt{A}, \texttt{B} et \texttt{n}).
  \item Donner le rang obtenu, puis montrer que le solde de B reste ensuite
  supérieur à celui de A en comparant les accroissements mensuels des deux offres.
\end{enumerate}

\end{exercice}


\noindent\textbf{Programme à exécuter et à analyser :}

\lstinputlisting[language=Python]{chapitres/1SPE-SUITES/python/1SPE-SUITES-TD-FIL-ROUGE-CROISEMENT.py}

\begin{corrige}{1SPE-SUITES-COURS-07-FR-EX3}

\textbf{Question 1.}
\[
  A_5 = 200 \times 5 = 1\,000 \text{ €}, \quad B_5 = 1\,000 \times 1{,}004^5 \approx 1\,020{,}16 \text{ €.}
\]
Donc $B_5 > A_5$ : grâce à son capital de départ, le solde de l'offre B est
encore le plus élevé au $5^{\text{e}}$ mois.
\[
  A_6 = 200 \times 6 = 1\,200 \text{ €}, \quad B_6 = 1\,000 \times 1{,}004^6 \approx 1\,024{,}24 \text{ €.}
\]
Donc $B_6 \leqslant A_6$ : dès le $6^{\text{e}}$ mois, le solde de l'offre A
devient supérieur à celui de l'offre B.

De plus,
\[
B_7\approx1\,028{,}34<A_7=1\,400,\quad
B_8\approx1\,032{,}45<A_8=1\,600,\quad
B_9\approx1\,036{,}58<A_9=1\,800.
\]
Ainsi, le solde de A reste supérieur aux rangs $6$ à $9$ ; la recherche qui
commence au rang $10$ ne laisse aucun croisement intermédiaire de côté.

\commentaireMarge{Le solde de B, supérieur au départ, est dépassé rapidement.}

On en conclut que le solde de B, supérieur au démarrage, est dépassé dès le
sixième mois. Il faut étudier des rangs beaucoup plus élevés pour savoir si
$B_n$ redevient supérieur à $A_n$.

\textbf{Question 2.} On sait par la question 1 que le solde de l'offre A est
supérieur à celui de B au rang $6$ et les calculs précédents couvrent les rangs $7$ à $9$ ; on initialise
au rang $10$ afin de chercher le croisement beaucoup
plus tardif. La boucle \texttt{while B <= A} tourne tant que le solde de l'offre B
n'est pas redevenu supérieur à celui de A. À chaque itération, on incrémente $n$ puis on recalcule les
deux capitaux avec les formules explicites. La boucle s'arrête au premier rang
$n\geqslant10$ tel que $B_n>A_n$.

\textbf{Question 3.} Exécution des deux premiers tours (après l'initialisation au rang $10$) :

\begin{center}
\begin{tabular}{|c|c|c|c|}
\hline
\texttt{n} & \texttt{A} & \texttt{B} & Condition \texttt{B <= A} \\
\hline
$10$ (init.) & $2\,000$ & $1\,040{,}73$ & Vrai (on entre dans la boucle) \\
\hline
$11$ & $2\,200$ & $1\,044{,}89$ & Vrai \\
\hline
$12$ & $2\,400$ & $1\,049{,}07$ & Vrai \\
\hline
\end{tabular}
\end{center}

\textbf{Question 4.} Le programme affiche \texttt{n = 1415}. Le premier rang
$n\geqslant10$ pour lequel $B_n>A_n$ est donc $n=1\,415$, soit environ
$118$ ans — un horizon qui dépasse très largement une vie humaine.

\[
  A_{1415} = 200 \times 1\,415 = 283\,000 \text{ €}, \quad B_{1415} = 1\,000 \times 1{,}004^{1415} \approx 283\,925 \text{ €.}
\]

Pour tout $n\geqslant1\,415$,
\[
B_{n+1}-B_n=0{,}004B_n
\quad\text{et}\quad
A_{n+1}-A_n=200.
\]
Comme $(B_n)$ est croissante et
$0{,}004B_{1415}\approx1\,135{,}70>200$, l'accroissement de l'offre B reste
supérieur à celui de l'offre A à partir de ce rang. Puisque
$B_{1415}>A_{1415}$, un raisonnement par récurrence montre alors que
$B_n>A_n$ pour tout $n\geqslant1\,415$.

\end{corrige}

\bigskip

%---------------------------------------------------------------
% PARTIE D — Synthese
%---------------------------------------------------------------

\begin{exercice}{1SPE-SUITES-COURS-07-FR-EX4}{3}{15}

\textbf{Partie D — Synthèse et interprétation économique}

\begin{enumerate}
  \item Quel livret présente le solde le plus élevé au bout de $5$ ans ? De
  $10$ ans ? Justifier avec les calculs des parties précédentes et rappeler que
  les deux modèles reposent sur des flux de versements différents.
  \item Expliquer pourquoi la comparaison des accroissements établie dans la
  partie C garantit que le solde de B reste supérieur après le rang $1\,415$.
  \item \pictoOral{} \textit{Question orale.} « Expliquer en quelques phrases
  les deux croisements observés entre les soldes. Quel rôle jouent le capital
  initial, le versement mensuel et les intérêts composés ? »
\end{enumerate}

\end{exercice}

\begin{corrige}{1SPE-SUITES-COURS-07-FR-EX4}

\textbf{Question 1.}

Sur $5$ ans ($n = 60$) : $A_{60} = 12\,000$~€ contre
$B_{60} \approx 1\,271$~€. Le solde de A est le plus élevé.

Sur $10$ ans ($n = 120$) : $A_{120} = 24\,000$~€ contre
$B_{120} = 1\,000 \times 1{,}004^{120} \approx 1\,615$~€. Le solde de A reste
le plus élevé.

\commentaireMarge{Les soldes comparent ici des flux différents : versements mensuels pour A, capital initial unique pour B.}

\textbf{Conclusion :} Sur les horizons de $5$ et $10$ ans étudiés, le solde de
A est supérieur. Ce constat décrit les deux flux choisis ; il ne compare pas
deux placements alimentés de la même manière. Les calculs couvrent les rangs
$6$ à $9$, puis le programme vérifie $A_n\geqslant B_n$ pour les rangs testés
de $10$ à $1\,414$.

\textbf{Question 2.} Au rang $1\,415$, le solde de B est supérieur. À partir de ce
rang, son accroissement mensuel est supérieur à $200$~€, tandis que celui de
l'offre A vaut exactement $200$~€. La preuve de la partie C montre donc que
l'écart en faveur de B ne peut plus se refermer.

\textbf{Question 3 — Éléments de réponse orale.}

Une suite géométrique de raison $q$ légèrement supérieure à $1$ croît lentement au début : les premiers termes ne font que multiplier le capital initial par des facteurs proches de $1$. En revanche, une suite arithmétique de grande raison $r$ accumule rapidement des montants importants par additions successives.

Le capital initial de l'offre B ($1\,000$~€) est modeste par rapport aux versements mensuels de l'offre A ($200$~€/mois). Pendant des années, les intérêts composés de $0{,}4\,\%$ par mois sur $1\,000$~€ restent inférieurs aux $200$~€ mensuels ajoutés dans l'offre A.

Au croisement observé au rang $1\,415$, les intérêts mensuels de l'offre B
dépassent enfin le versement fixe de $200$~€. La comparaison des accroissements
explique alors, dans ce cas précis, pourquoi l'offre B conserve son avance.

\end{corrige}
